\subsection{Vector functors}

In this number and in the three following numbers 7.7 to 7.9, the letter I denotes a finite set, the union of two disjoint subsets \(I_+\) and \(I_-\). We denote by \(\mathcal{V} = (\mathrm{V}_{i})_{i\in I}\) (and analogously by \(\mathcal{V}'\) , \(\mathcal{V}''\) , \ldots) a family of Banach spaces indexed by I. We denote by \(\operatorname{Hom}(\mathcal{V},\mathcal{V}')\) the Banach space \(\prod_{i\in I_+}\mathcal{L}(\mathrm{V}_i;\mathrm{V}_i')\times \prod_{i\in I_-}\mathcal{L}(\mathrm{V}_i';\mathrm{V}_i)\). We denote by \(\mathbf{f} = (f_i)\) an element of \(\operatorname{Hom}(\mathcal{V},\mathcal{V}')\). We write \(\mathrm{Id}_{\mathcal{V}}\) for the element \((\mathrm{Id}_{\mathrm{V}_i})_{i\in I}\) of \(\operatorname{Hom}(\mathcal{V},\mathcal{V})\). For \(\mathbf{f}\in \operatorname{Hom}(\mathcal{V},\mathcal{V}')\) and \(\mathbf{f}'\in \operatorname{Hom}(\mathcal{V}',\mathcal{V}'')\) , we denote by \(\mathbf{f}'\circ \mathbf{f}\) the element of \(\operatorname{Hom}(\mathcal{V},\mathcal{V}'')\) whose components are given by:

\[
\begin{array}{l l} (\mathbf {f} ^ {\prime} \circ \mathbf {f}) _ {i} = f _ {i} ^ {\prime} \circ f _ {i} & \text { if } \quad i \in \mathrm{I} _ {+} \\ (\mathbf {f} ^ {\prime} \circ \mathbf {f}) _ {i} = f _ {i} \circ f _ {i} ^ {\prime} & \text { if } \quad i \in \mathrm{I} _ {-}. \end{array}
\]

7.6.1. One calls a vector functor (resp. finite-dimensional vector functor) of type I and of class \(\mathbf{C}^r\) the data consisting, for every family \(\mathcal{V} = (\mathrm{V}_i)_{i\in I}\) of Banach spaces (resp. of finite-dimensional vector spaces over K), of a Banach space \(\tau(\mathcal{V})\) and, for every \(\mathbf{f}\in\mathrm{Hom}(\mathcal{V},\mathcal{V}')\), of an element \(\tau(\mathbf{f})\in\mathcal{L}(\tau(\mathcal{V});\tau(\mathcal{V}'))\) , these data being subject to the following two conditions:

\begin{enumerate}
\def\labelenumi{(\alph{enumi})}
\item
  One has \(\tau (\mathrm{Id}_{\mathcal{V}}) = \mathrm{Id}_{\tau (\mathcal{V})}\) and \(\tau (\mathbf{f}'\circ \mathbf{f}) = \tau (\mathbf{f}')\circ \tau (\mathbf{f}).\)
\item
  The map \(\tau : \operatorname{Hom}(\mathcal{V}, \mathcal{V}') \to \mathcal{L}(\tau(\mathcal{V}); \tau(\mathcal{V}'))\) is of class \(C'\).
\end{enumerate}

7.6.2. Let \(\mathcal{M} = (\mathbf{M}^i)_{i\in I}\) be a family of vector bundles with base B. For \(b\in \mathbf{B}\) , set \(\mathcal{M}_b = (\mathbf{M}_b^i)_{i\in I}\) . Let \(\tau\) be a vector functor and denote by \(\tau (\mathcal{M})\) the disjoint union of the \(\tau(\mathcal{M}_{b})\) for \(b\in\mathbf{B}\).

There exists on \(\tau(\mathcal{M})\) one vector bundle structure and only one (with base B relative to the map \(\pi\) of \(\tau(\mathcal{M})\) into B) such that for every \(b\in\mathbf{B}\) , one has \(\tau(\mathcal{M}_{b})=\pi^{-1}(b)\), possessing the following property:

Let U be an open subset of B and, for each i, let \(t^{i} = (\mathrm{U}, \varphi_{i}, \mathrm{F}_{i})\) be a vector chart of \(M^{i}\) with domain U; set \(\mathcal{F} = (\mathrm{F}_{i})_{i \in I}\) and let \(\psi_{b}\) be the element of \(\operatorname{Hom}(\mathcal{M}_{b}, \mathcal{F})\) defined by \((\psi_{b})_{i} = (t_{b}^{i})^{-1}\) for \(i \in I_{+}\) and \((\psi_{b})_{i} = t_{b}^{i}\) for \(i \in I_{-}\) ; for \(x \in \pi^{-1}(\mathrm{U})\) , set \(\psi(x) = (\pi(x), \tau(\psi_{\pi(x)})(x))\) . Then the triple \((\mathrm{U}, \psi, \tau(\mathcal{F}))\) is a vector chart of the vector bundle \(\tau(\mathcal{M})\) .

Equipped with this structure, \(\tau(\mathcal{M})\) is called the vector bundle deduced from the family \(\mathcal{M}\) by the vector functor \(\tau\) .

7.6.3. Let \(f\) be a morphism of \(\mathbf{B}\) into a manifold \(\mathbf{B}'\) . Let \(\mathcal{M} = (\mathbf{M}^i)\) (resp. \(\mathcal{M}' = (\mathbf{M}'^i)\) ) be a family indexed by \(\mathbf{I}\) of vector bundles over the base \(\mathbf{B}\) (resp. \(\mathbf{B}'\) ). For every \(i \in \mathbf{I}_+\) let \(g_i\) be an \(f\)-morphism of \(\mathbf{M}^i\) to \(\mathbf{M}'^i\) and for every \(i \in \mathbf{I}_-\) let \(g_i\) be an \(f\)-comorphism of \(\mathbf{M}'^i\) to \(\mathbf{M}^i\) . Set \(\mathbf{g} = (g_i)_{i \in I}\) and for \(b \in B\) , set \(\mathbf{g}_b = ((g_i)_b)_{i \in I}\) (cf. 7.2.1 and 7.2.6). There exists an \(f\)-morphism, and only one, denoted \(\tau(\mathbf{g})\) , from \(\tau(\mathcal{M})\) to \(\tau(\mathcal{M}')\) such that \(\tau(\mathbf{g})_b = \tau(\mathbf{g}_b)\) for all \(b \in B\) .

If in particular \(M^{i}=f^{*}M^{'i}\) , the \(g_{i}\) being the canonical morphisms or comorphisms, then the B-morphism from \(\tau(\mathcal{M})\) to \(f^{*}\tau(\mathcal{M}')\) defined by \(\tau(\mathbf{g})\) (7.2.4) is an isomorphism: one expresses this fact by saying that \(\tau\) commutes with inverse images.

In particular, let B' be a submanifold of B, and set \(\mathcal{M}|B' = (M^{i}|B')_{i\in I}\) . The vector bundles \(\tau(\mathcal{M})|B'\) and \(\tau(\mathcal{M}|B')\) are then canonically B'-isomorphic.

7.6.4. Let \(\tau, \tau_1, \ldots, \tau_d\) be vector functors (of type I and class \(C'\) ). A \(d\)-linear morphism \(\theta\) of \((\tau_1, \ldots, \tau_d)\) to \(\tau\) is the data consisting, for every family \(\mathcal{V}\) of Banach spaces indexed by I, of a continuous \(d\)-linear map \(\theta_{\mathcal{V}}\) from \(\tau_1(\mathcal{V}) \times \cdots \times \tau_d(\mathcal{V})\) to \(\tau(\mathcal{V})\), this data satisfying the following condition: for every \(f \in \text{Hom}(\mathcal{V}, \mathcal{V}')\) one has

\[
\tau (\mathbf {f}) \circ \theta_ {\mathcal{V}} = \theta_ {\mathcal{V}^{\prime}} \circ (\tau_ {1} (\mathbf {f}) \times \dots \times \tau _{d} (\mathbf {f})).
\]

For \(d = 1\) , one simply says a morphism of \(\tau_1\) to \(\tau\) .

Now let \(\mathcal{M}\) be a family indexed by I of vector bundles with base B. There exists one and only one B-morphism \(\theta_{\mathcal{M}}\), which is d-linear, from \(\tau_{1}(\mathcal{M}) \times_{\mathbf{B}} \cdots \times_{\mathbf{B}} \tau_{d}(\mathcal{M})\) to \(\tau(\mathcal{M})\) such that \((\theta_{\mathcal{M}})_{b} = \theta_{\mathcal{M}_{b}}\) for all \(b \in \mathbf{B}\) .

With the notations of 7.6.3., we have

\[
\tau (\mathbf {g}) \circ \theta_ {\mathcal {M}} = \theta_ {\mathcal {M} ^ {\prime}} \circ (\tau_ {1} (\mathbf {g}) \times \dots \times \tau_ {d} (\mathbf {g})).
\]

If \(d=1\) and if \(\theta\) is an isomorphism (which means that \(\theta_{\mathcal{V}}\) is an isomorphism for every family \(\mathcal{V}\) ), then \(\theta_{\mathcal{M}}\) is an isomorphism.

7.6.5. The definitions and results of nos. 7.6.2 to 7.6.4 extend to the case of vector functors in finite dimension, provided that the given vector bundles are everywhere assumed to be of finite rank.

They also extend to the following case: let L be a field equipped with a K-algebra structure of finite dimension; we take as \(\tau\) a vector functor over L (i.e., satisfying the hypotheses of 7.6.1 with K replaced by L), and we consider only vector bundles over L in the sense of 7.3.4.

7.6.6. A functor is called a vector functor (resp. finite-dimensional) for isomorphisms if, for every Banach space V (resp. every finite-dimensional vector space over K), it assigns a Banach space \(\tau(\mathrm{V})\) and for every isomorphism \(f\) from V onto a Banach space V' an isomorphism \(\tau(f)\) from \(\tau(\mathrm{V})\) onto \(\tau(\mathrm{V}')\) , these data being subject to condition (a) of 7.6.1 and to the following condition:

(b') The map \(f\mapsto\tau(f)\) from the open subset of \(\mathcal{L}(\mathrm{V};\mathrm{V}')\) consisting of the isomorphisms from V onto V' into \(\mathcal{L}(\tau(\mathrm{V});\tau(\mathrm{V}'))\) is of class C'.

The definitions and results of the preceding numbers extend to the case of vector functors for isomorphisms (by making \(I_{+} = \{1\}\) and \(I_{-} = \varnothing\) ), except for those in the first paragraph of no. 7.6.3.

